% !PRINTED 453
\weberpage{453}

If the $a$ are not independent variables, but are themselves in turn functions of other variables, all this remains valid, so long only as the changes of the $a$ take place continuously.

Accordingly, in all these cases we may speak of the $m$ roots of an equation of degree $m$, which are to be understood as numbers or as algebraic functions of their coefficients.

Now if $\Omega$ is an arbitrary field and $f(x)$ is a function in $\Omega$, two cases are possible.

The function $f(x)$ is either decomposable or not decomposable into factors of lower degree whose coefficients belong to the field $\Omega$. In the first case the function $f(x)$ is called \gesperrt{reducible}, in the second \gesperrt{irreducible} in $\Omega$. When the field $\Omega$ may be regarded as known and fixed, for instance when it is the field of rational numbers, one sometimes also speaks simply of reducible or irreducible functions, without expressly adding the designation ``in the field $\Omega$''. That in itself, however, the precise specification of the field to which the reducibility refers, or which is regarded as the domain of rationality, is very essential, follows from the following remarks.

A linear function is of course irreducible in every field. If one multiplies several functions of the form (1) with one another, there arises a function of the same form, which, however, is then naturally reducible, since it can again be decomposed into the factors from which it arose.

A function $f(x)$ which is irreducible in $\Omega$ can become reducible in an extended field $\Omega'$ which arises from $\Omega$ by any adjunction; thus every function $f(x)$ becomes reducible in the field $\Omega'$ which arises from $\Omega$ by the adjunction of a root $\alpha$ of $f(x) = 0$; for a linear factor $x - \alpha$ can be separated off from $f(x)$. In the field; which consists of \gesperrt{all} numbers, every function with numerical coefficients is reducible. % [sic] semicolon as printed

I. \gesperrt{An irreducible function $f(x)$ can have no common divisor with another function $F(x)$ whose coefficients belong to the same field, unless $F(x)$ is divisible by $f(x)$.}
